\Legend
先暂时固定所有棋子的初始点。因为每次只移动一枚棋子，而且图是 DAG，棋子游戏是 $k$ 个互不干扰的公平组合游戏之和。定义点 $v$ 的 Sprague–Grundy 值为其所有后继 SG 值集合的最小未出现非负整数，记为 $g(v)$。无棋可动时 SG 值为 $0$。固定初始点后的先手胜负正好由 $g(a_1)\oplus\cdots\oplus g(a_k)$ 是否非零决定。用拓扑序逆序求出全部 $g(v)$。

接着从最后一枚棋子向前处理位选择。设 $F_i(x)$ 表示前 $i-1$ 枚棋子的 SG 异或和为 $x$ 时，从第 $i$ 枚棋子的第 $0$ 位开始选择，小常最终能否获胜。边界为 $F_{k+1}(x)=[x\ne0]$。

处理一枚棋子时，先对每个可能点 $v$ 创建叶子，其值是按 $x$ 编号的布尔向量 $F_{i+1}(x\oplus g(v))$。然后按位从最高位到最低位倒推。若该位由小常决定，把选择 $0$ 和 $1$ 的两个向量按位或；若由小鹏决定，按位与。这样最低位最后合并，恰与真实的选择顺序一致。处理完 $k$ 枚棋子后答案为 $F_1(0)$。

为什么向量最多只需 $512$ 位？若某点 SG 值为 $G$，它至少有 $G$ 条出边通向 SG 值分别为 $0,1,\ldots,G-1$ 的点。这些较小值也各需要至少相应数量的出边，所以 $m\ge G(G+1)/2$。由 $m\le80000$ 得 $G\le399<512$，所有 SG 值的异或也小于 $512$。

把向量压成 $8$ 个 64 位字。每处理一枚棋子时，可先按不同的 SG 值缓存叶子的异或置换向量，再复制到各个点。复杂度为 $O(n+m+k(512^2+2^n\cdot 8))$，空间 $O(2^n\cdot8+512\cdot8+n+m)$。这里写成 $512^2$ 是对最多 $512$ 种 SG 值、每种扫描 $512$ 个异或状态的保守上界。
